\begin{lemma}[A valid three-edge switching preserves prescribed entries]
Let $G$ be a graph on a finite vertex set, with disjoint parts $U,W$.
Let $f:U\to W$ be a perfect matching as in \noderef{PerfectMatching}.
Take pairwise distinct $x,a,b\in U$ and $y\in W$ with $f(x)=y$.
Suppose $xf(a)$, $af(b)$, and $bf(x)$ are edges of $G$.
Let $(x_i,y_i)_{i\in I}$ be any family of prescribed entries such that
$f(x_i)=y_i$ and $x_i\notin\{x,a,b\}$ for every $i$.
Then $g=S(f;x,a,b)$, defined in \noderef{MatchingThreeEdgeSwitch}, is a
perfect matching, satisfies $g(x_i)=y_i$ for every $i$, and does not
contain the edge $xy$, that is, $g(x)\ne y$.
\end{lemma}
\begin{proof}
The matching definition in \noderef{PerfectMatching} says that $f$ is
bijective and each $zf(z)$ is an edge. By
\noderef{MatchingThreeEdgeSwitchProperties}, $g$ is bijective.
At $x,a,b$, its matching edges are respectively $xf(a)$, $af(b)$,
and $bf(x)$, which are edges by assumption. At every other $z$ its edge
is the original edge $zf(z)$. Hence it is a perfect matching.
Each $x_i$ is outside the three modified vertices, so the same value
identities give $g(x_i)=f(x_i)=y_i$. Finally $g(x)=f(a)$. If this were
$y=f(x)$, injectivity of $f$ would imply $a=x$, contradicting the
assumed distinctness. Thus $xy$ is removed.
\end{proof}
